\documentclass[a4paper,11pt]{article}

\usepackage[utf8]{inputenc}
\usepackage[T1]{fontenc}
\usepackage[ngerman]{babel}
\usepackage{graphicx}
\usepackage{array}
\usepackage{booktabs}
\usepackage{float}

\usepackage[scaled]{helvet}
\renewcommand{\familydefault}{\sfdefault}

\usepackage[a4paper, top=2cm, bottom=2cm, left=2cm, right=2cm]{geometry}

\usepackage{setspace}
\setstretch{1.5}

\setlength{\parindent}{0pt}
\setlength{\parskip}{6pt}

\usepackage{microtype}
\sloppy
\hyphenpenalty=1000
\tolerance=3000

\renewcommand{\footnotesize}{\fontsize{10}{12}\selectfont}

\setcounter{secnumdepth}{3}
\setcounter{tocdepth}{3}

\usepackage{titlesec}
\titleformat{\section}{\normalfont\fontsize{12}{14}\bfseries}{\thesection}{1em}{}
\titleformat{\subsection}{\normalfont\fontsize{12}{14}\bfseries}{\thesubsection}{1em}{}
\titleformat{\subsubsection}{\normalfont\fontsize{12}{14}\bfseries}{\thesubsubsection}{1em}{}

\usepackage[
  colorlinks=true,
  linkcolor=black,
  citecolor=blue,
  filecolor=black,
  urlcolor=blue
]{hyperref}
\usepackage[capitalise,nameinlink]{cleveref}

\usepackage{fancyhdr}
\pagestyle{fancy}
\fancyhf{}
\renewcommand{\headrulewidth}{0pt}
\fancyfoot[C]{\thepage}

\usepackage[backend=biber, style=apa]{biblatex}
\addbibresource{references.bib}

\usepackage{titling}

\usepackage{acronym}
\usepackage[german=quotes]{csquotes}

\usepackage{caption}
\captionsetup[table]{
    font=small,
    skip=10pt,
    labelfont=bf
}

\usepackage[table]{xcolor}
\definecolor{zebragray}{gray}{0.92}

\begin{document}

\begin{titlepage}
    \thispagestyle{empty}
    \centering
    \vspace*{5cm}
    {\Huge\bfseries Portfolio \par}
    \vspace{0.3cm}
    {\Large Konzeptionsphase (Phase~1) \par}
    \vspace{1cm}
    {\large Konzeptgrundlage für ein KMU-KI-Beratungsangebot: \\ Branchen, Use Cases und Geschäftsmodell \par}
    \vspace{1.5cm}
    {\large Kurs: DLBPKIEKPT01 -- Project: AI Fluency -- Einführung in die Generative KI (Code) \par}
    \vspace{0.5cm}
    {\large Studiengang: Angewandte Künstliche Intelligenz \par}
    \vspace{0.5cm}
    {\large Sven Behrens \par}
    \vspace{0.5cm}
    {\large Matrikelnummer: 42303511 \par}
    \vspace{0.5cm}
    {\large Tutor: [wird ergänzt] \par}
    \vspace{0.5cm}
    {\large \today \par}
\end{titlepage}

\pagenumbering{Roman}
\setcounter{page}{1}

\tableofcontents
\newpage

\pagenumbering{arabic}
\setcounter{page}{1}

% =============================================================================
% PHASE 1 -- KONZEPTIONSPHASE
% Laut Aufgabenstellung besteht Phase 1 aus drei Teilen:
%   1. Mein Problem            (ca. 1   Seite DIN A4)
%   2. Meine Strategie         (ca. 1-2 Seiten DIN A4)
%   3. Meine Tool-Auswahl      (ca. 0,5 Seiten DIN A4)
% =============================================================================

\section{Mein Problem}
% Anforderung: konkrete, persönlich relevante, inhaltlich überschaubare Aufgabe.

\subsection{Das Problem}
Mit Mitgründern befinde ich mich in der Frühphase eines
Gründungsvorhabens: Wir wollen ein Beratungsunternehmen aufbauen, das
kleinen lokalen Unternehmen (KMU) und Handwerksbetrieben hilft, generative
KI in konkrete Arbeitsabläufe einzubauen. Bevor wir uns auf eine
Branchenausrichtung, ein Leistungsangebot und ein Preismodell festlegen
können, müssen drei zusammenhängende Fragen beantwortet sein:

\begin{enumerate}
    \item Welche drei bis fünf lokalen Branchen kommen als realistische
    Erstkunden in Frage?
    \item Welche generativen KI-Use-Cases sind in diesen Branchen heute
    schon produktionsreif -- also nicht nur theoretisch denkbar, sondern
    mit überschaubarem Aufwand und akzeptablem Risiko produktiv
    einsetzbar?
    \item Wie übersetzt sich die Antwort auf (1) und (2) in unser eigenes
    Geschäftsmodell, dargestellt als Business Model Canvas?
\end{enumerate}

Im Rahmen dieses Portfolios bearbeite ich ausschließlich meinen Anteil an
dieser Konzeptarbeit, also die Vorbereitung des Canvas-Entwurfs samt der
zugrunde liegenden Branchen- und Use-Case-Recherche. Verhandlungen mit
Mitgründern, juristische Prüfungen, Behördengänge und der Aufbau der
operativen Struktur sind nicht Gegenstand dieser Arbeit; sie liegen
außerhalb dessen, was sich mit KI-Unterstützung sinnvoll innerhalb eines
Semesters bearbeiten lässt.

\subsection{Relevanz}
Das Problem ist für mich aus drei Gründen unmittelbar relevant. Erstens hat
unser Gründungsvorhaben reale Konsequenzen: Die hier entwickelte
Konzeptgrundlage fließt direkt in unsere ersten Akquisegespräche und in
unser Pitch-Material ein. Zweitens entsteht eine besondere
Authentizitätsfrage. Da wir KI-Beratungsleistungen verkaufen wollen,
müssen wir nachweislich selbst KI-kompetent sein -- ein ungeprüfter,
generisch übernommener KI-Output als Geschäftsgrundlage wäre inkonsistent
mit unserem eigenen Wertversprechen gegenüber zukünftigen Kunden.
Drittens fehlt mir aktuell genau das Methodengerüst, das die kritische,
urteilsfähige KI-Nutzung strukturiert, die für eine seriöse Konzeptarbeit
notwendig ist; das Modul \emph{Project: AI Fluency} bietet hierfür den
passenden Rahmen.

\subsection{Ziel und Abgrenzung}
Das Endergebnis dieser Portfolio-Arbeit besteht aus zwei Dokumenten:
einem ausgearbeiteten Business Model Canvas für unser geplantes
KMU-KI-Beratungsangebot sowie einem Branchen- und Use-Case-Mapping als
Anlage, das die Auswahl der Erstzielbranchen und der vorgeschlagenen
Anwendungsfälle nachvollziehbar begründet.

Ein gutes Ergebnis erfüllt aus meiner Sicht drei Kriterien. Erstens ist
die Branchenauswahl mit eigenem regionalen Kontextwissen und eigener
Recherche begründet, nicht allein mit KI-generierten Vorschlägen.
Zweitens unterscheiden die genannten Use-Cases klar zwischen
produktionsreifen, eingeschränkt einsatzfähigen und (noch) ungeeigneten
Anwendungen, inklusive einer offenen Beschreibung, womit KI heute
scheitert. Drittens ist der Canvas eine begründete Arbeitshypothese, die
in nachfolgenden Validierungsgesprächen geprüft werden kann -- kein
fertiges Geschäftsmodell. Diese Zurückhaltung im Anspruch ist Teil der
inhaltlichen Aussage: Eine Konzeptphase liefert keinen validierten Markt,
sondern eine prüfbare Arbeitsgrundlage.

\section{Meine Strategie: Mensch-KI-Aufgabenteilung}
% Anforderung: Plane BEVOR Du loslegst.
% Drei Tabellen abdecken:
%   (a) Was ICH selbst mache       -> Warum menschliche Kompetenz?
%   (b) Wo KI mich unterstützt     -> Warum ist KI hier sinnvoll?
%   (c) Wo ich KEINE KI nutze      -> Warum nicht?

[Einleitende Argumentation folgt.]

\begin{table}[H]
\centering
\caption{Was ich selbst mache}
\label{tab:strategie_mensch}
\rowcolors{2}{zebragray}{white}
\begin{tabular}{p{6cm}p{9cm}}
\toprule
\textbf{Tätigkeit} & \textbf{Warum menschliche Kompetenz nötig ist} \\
\midrule
\textit{z.\,B.\ Zieldefinition}            & \textit{Begründung} \\
\textit{z.\,B.\ Strategische Entscheidung} & \textit{Begründung} \\
\textit{z.\,B.\ Qualitätskontrolle}        & \textit{Begründung} \\
\bottomrule
\end{tabular}
\end{table}

\begin{table}[H]
\centering
\caption{Wo KI mich unterstützt}
\label{tab:strategie_ki}
\rowcolors{2}{zebragray}{white}
\begin{tabular}{p{6cm}p{9cm}}
\toprule
\textbf{Tätigkeit} & \textbf{Warum KI hier sinnvoll ist} \\
\midrule
\textit{z.\,B.\ Recherche}        & \textit{Begründung} \\
\textit{z.\,B.\ Erste Entwürfe}   & \textit{Begründung} \\
\textit{z.\,B.\ Strukturierung}   & \textit{Begründung} \\
\bottomrule
\end{tabular}
\end{table}

\begin{table}[H]
\centering
\caption{Wo ich KEINE KI nutze}
\label{tab:strategie_ohneki}
\rowcolors{2}{zebragray}{white}
\begin{tabular}{p{6cm}p{9cm}}
\toprule
\textbf{Tätigkeit} & \textbf{Warum nicht} \\
\midrule
\textit{z.\,B.\ Finale Entscheidungen} & \textit{Begründung} \\
\textit{z.\,B.\ Ethische Abwägungen}   & \textit{Begründung} \\
\bottomrule
\end{tabular}
\end{table}

[Abschließende Begründung folgt.]

\section{Meine Tool-Auswahl}
% Anforderung: Welche Tools nutzt Du -- und warum?
%   - KI-Tools (welche und warum)
%   - Andere Ressourcen (Fachliteratur, Experten, etc.)
% Format frei: Text, Tabelle, Mindmap.

[Inhalt folgt: begründete Auswahl der eingesetzten KI-Tools sowie weiterer
Ressourcen (z.\,B.\ Gespräche mit lokalen Unternehmern, Branchenliteratur).]

\newpage

\printbibliography
\addcontentsline{toc}{section}{Literaturverzeichnis}

\end{document}
